\chapter{حلقه‌ها}
\label{ch:loops}

\begin{goals}
\begin{itemize}
  \item تکرار یک کار به تعداد مشخص با \fcode{تکرار}
  \item شمارش از یک عدد تا عدد دیگر، رو به بالا، رو به پایین و با گام
  \item تکرار تا وقتی که شرطی برقرار است، با \fcode{تاوقتی}
  \item شکستن حلقه با \fcode{بشکن} و پریدن از یک دور با \fcode{ادامه}
  \item حلقه درون حلقه
\end{itemize}
\end{goals}

\section{چرا حلقه؟}

فرض کنید می‌خواهید عددهای ۱ تا ۵ را چاپ کنید. با آنچه تا حالا یاد
گرفته‌اید، پنج دستور \fcode{چاپ} می‌نویسید. حالا اگر بخواهید تا ۱۰۰۰ را چاپ
کنید چه؟ یا اگر تعداد از پیش معلوم نباشد؟ رایانه‌ها دقیقاً برای همین
کارهای تکراری ساخته شده‌اند، و \textbf{حلقه} ابزاری است که به آن‌ها
می‌گوید: «این کار را بارها انجام بده.»

\section{تکرار به تعداد مشخص}

ساده‌ترین حلقه‌ی سلام \kwd{تکرار} است. جلویش می‌نویسید چند بار، و سلام بلوک
را همان تعداد بار اجرا می‌کند:

\program{ساده‌ترین حلقه}
\begin{salam}
تابع اصلی:
    تکرار ۳:
        چاپ "سلام!"
    پایان
    چاپ "پایان"
پایان
\end{salam}

\begin{outputfa}
سلام!
سلام!
سلام!
پایان
\end{outputfa}

ساختار حلقه همان ساختار آشناست: خطی که با دونقطه تمام می‌شود، بلوک با
تورفتگی و \fcode{پایان}. به هر بار اجرای بلوک یک \textbf{دور} می‌گوییم. این
حلقه سه دور دارد. عدد جلوی \fcode{تکرار} می‌تواند یک متغیر یا محاسبه هم
باشد: \fcode{تکرار تعداد مهمان:} یا \fcode{تکرار ن * ۲:}.

\section{شمارنده‌ی حلقه}

بیشتر وقت‌ها می‌خواهیم بدانیم در کدام دور هستیم. با نوشتن \kwd{با} و یک
نام، سلام در هر دور شماره‌ی آن دور را در یک متغیر می‌گذارد. شماره از
\emph{صفر} شروع می‌شود:

\program{شمارنده‌ی حلقه}
\begin{salam}
تابع اصلی:
    تکرار ۵ با شماره:
        چاپ "دور", شماره
    پایان
پایان
\end{salam}

\begin{outputfa}
دور 0
دور 1
دور 2
دور 3
دور 4
\end{outputfa}

پنج دور، با شماره‌های ۰ تا ۴. شروع از صفر شاید عجیب باشد، اما در
برنامه‌نویسی یک رسم بسیار قدیمی است و در فصل \ref{ch:arrays} خواهید دید که
چرا مفید است. متغیر \fcode{شماره} فقط درون حلقه وجود دارد و بعد از
\fcode{پایان} دیگر در دسترس نیست.

\section{شمارش از یک عدد تا عدد دیگر}

اگر بخواهید از ۱ بشمارید نه از ۰، یا از ۱۰ تا ۲۰، از شکل \fcode{تکرار آغاز
تا پایان} استفاده کنید. هر دو سر بازه شمرده می‌شوند:

\program{شمارش از یک تا پنج}
\begin{salam}
تابع اصلی:
    تکرار ۱ تا ۵ با شماره:
        چاپ شماره, "ضرب در ۳ می‌شود", شماره * ۳
    پایان
پایان
\end{salam}

\begin{outputfa}
1 ضرب در ۳ می‌شود 3
2 ضرب در ۳ می‌شود 6
3 ضرب در ۳ می‌شود 9
4 ضرب در ۳ می‌شود 12
5 ضرب در ۳ می‌شود 15
\end{outputfa}

اگر عدد آغاز از پایان بزرگ‌تر باشد، سلام رو به پایین می‌شمارد:

\program{شمارش معکوس}
\begin{salam}
تابع اصلی:
    تکرار ۵ تا ۱ با ثانیه:
        چاپ ثانیه
    پایان
    چاپ "پرتاب!"
پایان
\end{salam}

\begin{outputfa}
5
4
3
2
1
پرتاب!
\end{outputfa}

و با \kwd{گام} می‌توان به جای یکی‌یکی، چندتاچندتا شمرد:

\program{شمارش با گام}
\begin{salam}
تابع اصلی:
    متغیر خط := ""
    تکرار ۰ تا ۲۰ گام ۵ با عدد:
        خط = خط + عدد + " "
    پایان
    چاپ خط
پایان
\end{salam}

\begin{outputfa}
0 5 10 15 20
\end{outputfa}

در این برنامه به جای چاپ هر عدد در یک خط، همه را در یک رشته کنار هم
چیدیم و در آخر یک بار چاپ کردیم. این ترفند (شروع با رشته‌ی خالی و چسباندن
قطعه‌ها) را بارها خواهید دید. گام باید عددی مثبت باشد؛ برای شمارش معکوس با
گام، کافی است آغاز را بزرگ‌تر از پایان بگذارید.

\section{جمع کردن در حلقه}

حالا الگوی «جمع کردن» فصل \ref{ch:variables} را با حلقه ترکیب کنیم. یک
متغیر تغییرپذیر با مقدار صفر می‌سازیم و در هر دور چیزی به آن اضافه
می‌کنیم:

\program{مجموع یک تا صد}
\begin{salam}
تابع اصلی:
    متغیر مجموع := ۰
    تکرار ۱ تا ۱۰۰ با عدد:
        مجموع += عدد
    پایان
    چاپ "مجموع ۱ تا ۱۰۰:", مجموع
پایان
\end{salam}

\begin{outputfa}
مجموع ۱ تا ۱۰۰: 5050
\end{outputfa}

دقت کنید که \fcode{مجموع} حتماً باید \fcode{متغیر} باشد، چون در هر دور
عوض می‌شود، و حتماً باید \emph{بیرون} حلقه تعریف شود؛ اگر درون حلقه
تعریفش می‌کردیم، در هر دور یک جعبه‌ی تازه با مقدار صفر ساخته می‌شد و
هیچ‌وقت چیزی جمع نمی‌شد.

حلقه و شرط با هم، ترکیب بسیار قدرتمندی می‌سازند. مثلاً جمع و شمارش فقط
عددهای زوج:

\program{شمارش و جمع زوج‌ها}
\begin{salam}
تابع اصلی:
    متغیر مجموع زوج := ۰
    متغیر تعداد زوج := ۰
    تکرار ۱ تا ۲۰ با عدد:
        اگر عدد % ۲ == ۰:
            مجموع زوج += عدد
            تعداد زوج++
        پایان
    پایان
    چاپ "تعداد:", تعداد زوج, "مجموع:", مجموع زوج
پایان
\end{salam}

\begin{outputfa}
تعداد: 10 مجموع: 110
\end{outputfa}

\section{تکرار تا وقتی که: تاوقتی}

همه‌ی حلقه‌های بالا از قبل می‌دانستند چند دور دارند. اما گاهی نمی‌دانیم:
«آن‌قدر پس‌انداز کن تا به یک میلیون برسی.» «آن‌قدر عدد را بر ده تقسیم کن تا
صفر شود.» برای این حلقه‌ها واژه‌ی \kwd{تاوقتی} هست. جلویش یک شرط می‌آید و
تا وقتی شرط درست است، بلوک تکرار می‌شود:

\program{پس‌انداز تا رسیدن به هدف}
\begin{salam}
تابع اصلی:
    متغیر پس انداز := ۱۰۰۰۰۰
    متغیر ماه := ۰
    تاوقتی پس انداز < ۱۰۰۰۰۰۰:
        پس انداز += ۱۵۰۰۰۰
        ماه++
    پایان
    چاپ "بعد از", ماه, "ماه پس‌انداز به", پس انداز, "رسید"
پایان
\end{salam}

\begin{outputfa}
بعد از 6 ماه پس‌انداز به 1000000 رسید
\end{outputfa}

سلام پیش از \emph{هر} دور شرط را می‌سنجد. تا وقتی پس‌انداز از یک میلیون
کمتر است، ۱۵۰ هزار به آن اضافه می‌کند و یک ماه می‌شمارد. به محض این که شرط
نادرست شد، حلقه تمام می‌شود و برنامه از بعد از \fcode{پایان} ادامه می‌دهد.
اگر شرط از همان اول نادرست باشد، بلوک حتی یک بار هم اجرا نمی‌شود.

\program{شمردن رقم‌های یک عدد}
\begin{salam}
تابع اصلی:
    عدد := ۴۸۲۱۳
    متغیر باقی := عدد
    متغیر تعداد رقم := ۰
    تاوقتی باقی > ۰:
        باقی /= ۱۰
        تعداد رقم++
    پایان
    چاپ عدد, "دارای", تعداد رقم, "رقم است"
پایان
\end{salam}

\begin{outputfa}
48213 دارای 5 رقم است
\end{outputfa}

هر تقسیم صحیح بر ۱۰، یک رقم از سمت راست عدد می‌اندازد: ۴۸۲۱۳، ۴۸۲۱، ۴۸۲،
۴۸، ۴، ۰. پنج بار طول کشید تا به صفر برسد، پس عدد پنج رقم دارد. متغیر
\fcode{باقی} را برای این ساختیم که خود \fcode{عدد} دست‌نخورده بماند و در
آخر بتوانیم چاپش کنیم.

\begin{warn}[حلقه‌ی بی‌پایان]
در حلقه‌ی \fcode{تاوقتی}، چیزی درون بلوک باید شرط را سرانجام نادرست کند.
اگر خط \fcode{باقی \lr{/=} ۱۰} را از برنامه‌ی بالا حذف کنید، \fcode{باقی}
هرگز صفر نمی‌شود و حلقه تا ابد می‌چرخد. به این وضعیت \textbf{حلقه‌ی
بی‌پایان} می‌گویند و یکی از رایج‌ترین اشتباه‌های تازه‌کارهاست. اگر در
ویرایشگر برخط برنامه‌ای «در حال اجرا» ماند و تمام نشد، احتمالاً به همین
مشکل برخورده‌اید؛ صفحه را دوباره بارگذاری کنید و به دنبال خطی بگردید که
باید شرط را عوض می‌کرده است.
\end{warn}

\begin{note}[کدام حلقه؟]
اگر تعداد دورها از قبل معلوم است (یا از روی یک عدد به دست می‌آید)،
\fcode{تکرار} انتخاب طبیعی است. اگر تعداد دورها به چیزی بستگی دارد که در
طول حلقه رخ می‌دهد، \fcode{تاوقتی} را برگزینید. هر کاری را با هر دو می‌شود
کرد، اما یکی از آن دو همیشه خواناتر است.
\end{note}

\section{شکستن حلقه و پریدن از یک دور}

گاهی در میانه‌ی حلقه می‌فهمیم که دیگر لازم نیست ادامه دهیم. دستور
\kwd{بشکن} حلقه را همان لحظه تمام می‌کند:

\program{پیدا کردن و متوقف شدن}
\begin{salam}
تابع اصلی:
    تکرار ۱ تا ۱۰ با عدد:
        اگر عدد * عدد > ۳۰:
            چاپ "اولین عددی که مربعش از ۳۰ بیشتر است:", عدد
            بشکن
        پایان
    پایان
پایان
\end{salam}

\begin{outputfa}
اولین عددی که مربعش از ۳۰ بیشتر است: 6
\end{outputfa}

بدون \fcode{بشکن}، حلقه تا ۱۰ ادامه می‌داد و برای ۷، ۸، ۹ و ۱۰ هم پیام
چاپ می‌کرد. \fcode{بشکن} گفت: «پیدا شد، بس است.»

دستور \kwd{ادامه} ملایم‌تر است: فقط از \emph{همین دور} می‌گذرد و به دور
بعد می‌رود:

\program{گذشتن از مضرب‌های سه}
\begin{salam}
تابع اصلی:
    تکرار ۱ تا ۱۰ با عدد:
        اگر عدد % ۳ == ۰: ادامه پایان
        بنویس عدد, ""
    پایان
    چاپ
پایان
\end{salam}

\begin{outputfa}
1 2 4 5 7 8 10
\end{outputfa}

عددهای ۳، ۶ و ۹ چاپ نشدند، چون برای آن‌ها \fcode{ادامه} اجرا شد و
\fcode{بنویس} هرگز به نوبت نرسید. (ترفند کوچک: \fcode{بنویس عدد, ""}
عدد را چاپ می‌کند و بعدش یک فاصله، چون میان چیزهایی که با ویرگول جدا
شده‌اند یک فاصله می‌گذارد. \fcode{چاپ} خالی در آخر هم خط را می‌بندد.)

\section{حلقه درون حلقه}

مثل شرط، حلقه هم می‌تواند درون حلقه‌ی دیگری باشد. حلقه‌ی درونی برای هر دور
حلقه‌ی بیرونی، \emph{کامل} اجرا می‌شود. کلاسیک‌ترین نمونه، جدول ضرب است:

\program{جدول ضرب}
\begin{salam}
تابع اصلی:
    تکرار ۱ تا ۵ با سطر:
        متغیر خط := ""
        تکرار ۱ تا ۵ با ستون:
            خط = خط + (سطر * ستون) + " "
        پایان
        چاپ خط
    پایان
پایان
\end{salam}

\begin{outputfa}
1 2 3 4 5
2 4 6 8 10
3 6 9 12 15
4 8 12 16 20
5 10 15 20 25
\end{outputfa}

حلقه‌ی بیرونی پنج دور دارد؛ در هر دور، حلقه‌ی درونی پنج بار می‌چرخد و یک
خط از جدول را می‌سازد. در کل ۲۵ ضرب انجام می‌شود. متغیر \fcode{خط} را
\emph{درون} حلقه‌ی بیرونی تعریف کردیم تا برای هر سطر، تازه و خالی باشد.

\program{مثلث ستاره‌ای}
\begin{salam}
تابع اصلی:
    تکرار ۱ تا ۴ با ردیف:
        متغیر خط := ""
        تکرار ردیف:
            خط += "*"
        پایان
        چاپ خط
    پایان
پایان
\end{salam}

\begin{outputfa}
*
**
***
****
\end{outputfa}

این‌جا تعداد دورهای حلقه‌ی درونی به شمارنده‌ی حلقه‌ی بیرونی بستگی دارد:
ردیف اول یک ستاره، ردیف دوم دو ستاره و همین‌طور تا آخر. با تغییرهای کوچک
در همین برنامه می‌توان شکل‌های گوناگونی ساخت؛ تمرین‌های پایان فصل چند
نمونه پیشنهاد می‌کنند.

\section{یک برنامه‌ی کامل‌تر}

جدول تبدیل دما از سلسیوس به فارنهایت، از صفر تا چهل درجه، ده تا ده تا:

\program{جدول تبدیل دما}
\begin{salam}
تابع اصلی:
    چاپ "سلسیوس -> فارنهایت"
    چاپ "-------------------"
    تکرار ۰ تا ۴۰ گام ۱۰ با سلسیوس:
        فارنهایت := (سلسیوس بعنوان اعشار۶۴) * ۹.۰ / ۵.۰ + ۳۲.۰
        چاپ سلسیوس, "->", فارنهایت
    پایان
پایان
\end{salam}

\begin{outputfa}
سلسیوس -> فارنهایت
-------------------
0 -> 32
10 -> 50
20 -> 68
30 -> 86
40 -> 104
\end{outputfa}

شمارنده‌ی حلقه عددی صحیح است، پس پیش از محاسبه‌ی اعشاری آن را
\fcode{بعنوان اعشار۶۴} تبدیل کردیم. متغیر \fcode{فارنهایت} درون حلقه و بدون
\fcode{متغیر} تعریف شده است؛ در هر دور جعبه‌ای تازه ساخته می‌شود، مقدار
می‌گیرد، چاپ می‌شود و با پایان دور از بین می‌رود.

\begin{summary}
\begin{itemize}
  \item \fcode{تکرار ن:} بلوک را ن بار اجرا می‌کند. با \fcode{با شماره}
        شماره‌ی دور (از صفر) را می‌گیرید.
  \item \fcode{تکرار آغاز تا پایان با شماره:} از آغاز تا پایان می‌شمارد،
        هر دو سر شامل؛ اگر آغاز بزرگ‌تر باشد، رو به پایین. \fcode{گام}
        اندازه‌ی پرش را مشخص می‌کند.
  \item \fcode{تاوقتی شرط:} تا وقتی شرط درست است تکرار می‌کند. چیزی درون
        حلقه باید شرط را سرانجام نادرست کند.
  \item \fcode{بشکن} حلقه را تمام می‌کند؛ \fcode{ادامه} به دور بعد می‌پرد.
  \item حلقه‌ها می‌توانند تودرتو باشند؛ حلقه‌ی درونی برای هر دور بیرونی
        کامل اجرا می‌شود.
  \item متغیری که در حلقه جمع می‌شود، باید \fcode{متغیر} و بیرون حلقه تعریف
        شود.
\end{itemize}
\end{summary}

\section*{تمرین‌ها}

\exercise عددهای ۱ تا ۵۰ را که بر ۷ بخش‌پذیرند، در یک خط چاپ کنید.

\exercise حاصل ضرب عددهای ۱ تا ۱۰ را با حلقه حساب و چاپ کنید.

\exercise مبلغی را در بانک با سود سالانه‌ی ۲۰ درصد می‌گذارید. با
\fcode{تاوقتی} حساب کنید بعد از چند سال مبلغ دو برابر می‌شود و در هر سال
موجودی را چاپ کنید.

\exercise برنامه‌ی مثلث ستاره‌ای را طوری تغییر دهید که مثلث وارونه بسازد
(اول چهار ستاره، بعد سه، \ldots). سپس نسخه‌ای بنویسید که یک مربع توخالی
۵ در ۵ از ستاره بسازد.

\exercise با دو حلقه‌ی تودرتو همه‌ی جفت‌عددهای میان ۱ تا ۱۰ را پیدا و چاپ
کنید که مجموعشان ۱۰ است (مثل ۱ و ۹، ۲ و ۸). هر جفت فقط یک بار چاپ شود.

\exercise برنامه‌ی زیر چند خط چاپ می‌کند؟ اول حساب کنید، بعد اجرا کنید:
\begin{salam}
تابع اصلی:
    تکرار ۳ با الف:
        تکرار ۴ با ب:
            اگر ب == ۲: بشکن پایان
            چاپ الف, ب
        پایان
    پایان
پایان
\end{salam}
